\ProvidesFile{estilo.tex}[Estilo reutilizable para informes académicos]

% Este archivo se carga desde la raíz del proyecto con \input.
\RequirePackage{fontspec}
\RequirePackage[spanish]{babel}
\RequirePackage[a4paper,margin=20mm,headheight=14pt,headsep=6mm,footskip=12mm]{geometry}
\RequirePackage{graphicx}
\RequirePackage{xcolor}
\RequirePackage{array}
\RequirePackage{longtable}
\RequirePackage{tabularx}
\RequirePackage{booktabs}
\RequirePackage{enumitem}
\RequirePackage{pdflscape}
\RequirePackage{listings}
\RequirePackage{tcolorbox}
\tcbuselibrary{breakable,listings}
\RequirePackage{xurl}
\RequirePackage{fancyhdr}
\RequirePackage{hyperref}

% No se permite sustituir Montserrat por una fuente del sistema.
\newcommand{\InformeCompruebaFuente}[1]{%
  \IfFileExists{recursos/fuentes/Montserrat/#1}{}{%
    \PackageError{estilo}{Falta la fuente local Montserrat: #1}%
      {Compila desde la raiz del proyecto y restaura los archivos de recursos/fuentes/Montserrat.}%
  }%
}
\InformeCompruebaFuente{Montserrat-Regular.ttf}
\InformeCompruebaFuente{Montserrat-Bold.ttf}
\InformeCompruebaFuente{Montserrat-Italic.ttf}
\InformeCompruebaFuente{Montserrat-BoldItalic.ttf}
\InformeCompruebaFuente{Montserrat-SemiBold.ttf}
\InformeCompruebaFuente{Montserrat-SemiBoldItalic.ttf}

\setmainfont{Montserrat}[
  Path = recursos/fuentes/Montserrat/,
  Extension = .ttf,
  UprightFont = Montserrat-Regular,
  BoldFont = Montserrat-Bold,
  ItalicFont = Montserrat-Italic,
  BoldItalicFont = Montserrat-BoldItalic
]
\newfontfamily\InformeSemibold{Montserrat}[
  Path = recursos/fuentes/Montserrat/,
  Extension = .ttf,
  UprightFont = Montserrat-SemiBold,
  ItalicFont = Montserrat-SemiBoldItalic
]
% Se conserva la familia monoespaciada de TeX para listings y comandos.

\definecolor{InformeNavy}{HTML}{19324A}
\definecolor{InformeTeal}{HTML}{167D86}
\definecolor{InformeTexto}{HTML}{27343B}
\definecolor{InformeGrisClaro}{HTML}{F2F5F7}
\definecolor{InformeBorde}{HTML}{CBD5DC}

\AtBeginDocument{\fontsize{10pt}{13pt}\selectfont}
\color{InformeTexto}
\setlength{\emergencystretch}{2em}
\setlength{\parindent}{0pt}
\setlength{\parskip}{4pt}
\setlength{\tabcolsep}{3.8pt}
\renewcommand{\arraystretch}{1.12}
\setlength{\LTpre}{4pt}
\setlength{\LTpost}{4pt}
\setlength{\LTleft}{0pt}
\setlength{\LTright}{0pt}
\setlist{itemsep=2pt,topsep=4pt,parsep=0pt,partopsep=0pt}

\newcolumntype{Y}{>{\raggedright\arraybackslash}X}
\newenvironment{TablaCompacta}{%
  \begingroup
  \small
  \setlength{\tabcolsep}{3.8pt}%
  \renewcommand{\arraystretch}{1.08}%
}{\endgroup}

\lstdefinestyle{InformeCodigo}{%
  language={},
  basicstyle=\ttfamily\footnotesize,
  breaklines=true,
  breakatwhitespace=false,
  columns=fullflexible,
  keepspaces=true,
  showstringspaces=false,
  upquote=true
}
\newtcblisting{CodigoTerminal}{%
  listing only,
  breakable,
  colback=InformeGrisClaro,
  colframe=InformeBorde,
  boxrule=0.4pt,
  arc=1mm,
  left=2mm,
  right=2mm,
  top=1mm,
  bottom=1mm,
  listing options={style=InformeCodigo}
}
\newtcolorbox{RecuadroInforme}[1][]{%
  breakable,
  colback=InformeGrisClaro,
  colframe=InformeBorde,
  coltitle=InformeNavy,
  fonttitle=\InformeSemibold,
  boxrule=0.4pt,
  arc=1mm,
  left=2mm,
  right=2mm,
  top=1.5mm,
  bottom=1.5mm,
  #1
}

% Portadilla parametrizada: institución, título, subtítulo, autoría, unidad y fecha.
\newcommand{\PortadaInforme}[6]{%
  \begin{titlepage}
    \thispagestyle{empty}
    \centering
    \vspace*{12mm}
    {\InformeSemibold\Large #1\par}
    \vfill
    {\color{InformeTeal}\rule{0.22\textwidth}{2pt}\par}
    \vspace{8mm}
    {\color{InformeNavy}\bfseries\fontsize{25}{30}\selectfont #2\par}
    \vspace{6mm}
    {\large #3\par}
    \vfill
    {\InformeSemibold #4\par}
    {\small #5\par}
    \vspace{10mm}
    {\small #6\par}
    \vspace*{10mm}
  \end{titlepage}%
}

% Las cuatro etiquetas quedan vacías hasta que el autor las configure.
\newcommand{\InformeEncabezadoIzquierdo}{}
\newcommand{\InformeEncabezadoDerecho}{}
\newcommand{\InformePieIzquierdo}{}
\newcommand{\InformePieDerecho}{}
\newcommand{\InformeConfiguraEncabezados}{%
  \pagestyle{fancy}
  \fancyhf{}
  \fancyhead[L]{\small\nouppercase{\InformeEncabezadoIzquierdo}}
  \fancyhead[R]{\small\nouppercase{\InformeEncabezadoDerecho}}
  \fancyfoot[L]{\small\nouppercase{\InformePieIzquierdo}}
  \fancyfoot[C]{\small\thepage}
  \fancyfoot[R]{\small\nouppercase{\InformePieDerecho}}
  \renewcommand{\headrulewidth}{0.4pt}
  \renewcommand{\footrulewidth}{0pt}
  \fancypagestyle{plain}{%
    \fancyhf{}
    \fancyhead[L]{\small\nouppercase{\InformeEncabezadoIzquierdo}}
    \fancyhead[R]{\small\nouppercase{\InformeEncabezadoDerecho}}
    \fancyfoot[L]{\small\nouppercase{\InformePieIzquierdo}}
    \fancyfoot[C]{\small\thepage}
    \fancyfoot[R]{\small\nouppercase{\InformePieDerecho}}
    \renewcommand{\headrulewidth}{0.4pt}
    \renewcommand{\footrulewidth}{0pt}
  }
}
\InformeConfiguraEncabezados

\newcommand{\FuenteFigura}[1]{%
  \par\smallskip{\footnotesize\color{InformeTexto}\textit{Fuente:} #1\par}%
}
\newcommand{\FiguraEvidencia}[4]{%
  \begin{figure}[tbp]
    \centering
    \includegraphics[width=\linewidth,height=0.72\textheight,keepaspectratio]{#1}
    \caption{#2}\label{#4}
    \FuenteFigura{#3}
  \end{figure}%
}

\newcommand{\RefFigura}[1]{Figura~\ref{#1}}
\newcommand{\RefTabla}[1]{Tabla~\ref{#1}}
\newcommand{\RefAnexo}[1]{Anexo~\ref{#1}}
\newcommand{\RefPagina}[1]{página~\pageref{#1}}
\newcommand{\IniciarAnexos}{\clearpage\appendix}
\newcommand{\Anexo}[1]{\section{#1}}
\newenvironment{PaginaApaisada}{\begin{landscape}}{\end{landscape}}

\hypersetup{
  colorlinks=true,
  hypertexnames=false,
  linkcolor=InformeNavy,
  citecolor=InformeTeal,
  urlcolor=InformeTeal,
  pdfcreator={XeLaTeX}
}

% Reserve room for the report's Roman section labels in the contents.
\makeatletter
\renewcommand*{\l@section}{\@dottedtocline{1}{1.5em}{3.6em}}
\renewcommand*{\l@subsection}{\@dottedtocline{2}{3.8em}{4.3em}}
\makeatother
